\chapter{pip}

\section{下载}
到\href{https://pypi.org/project/pip/}{pip官网}下载pip安装包。
\begin{lstlisting}[style=demostyle,numbers=none]
# 下载pip安装包
wget https://bootstrap.pypa.io/get-pip.py
# 安装pip
python get-pip.py

# 查看pip版本
pip --version # Python2.x 版本命令
pip3 --version    # Python3.x 版本命令

# 升级pip
## Linux 或 macOS
pip install --upgrade pip    # python2.x
pip3 install --upgrade pip   # python3.x

## Windows 平台升级：
python -m pip install -U pip   # python2.x
python -m pip3 install -U pip    # python3.x
\end{lstlisting}

\section{国内源}
\begin{lstlisting}[style=demostyle,numbers=none]
阿里云 http://mirrors.aliyun.com/pypi/simple/ 
中国科技大学 https://pypi.mirrors.ustc.edu.cn/simple/ 
豆瓣(douban) http://pypi.douban.com/simple/ 
清华大学 https://pypi.tuna.tsinghua.edu.cn/simple/ 
中国科学技术大学 http://pypi.mirrors.ustc.edu.cn/simple/
\end{lstlisting}
\subsection{配置国内源}
\begin{lstlisting}[style=demostyle,numbers=none]
# 临时生效
pip install -i https://pypi.tuna.tsinghua.edu.cn/simple SomePackage

# 永久生效
pip config set global.index-url https://pypi.tuna.tsinghua.edu.cn/simple
\end{lstlisting}

\begin{enumerate}[label=\arabic*)]
    \item Linux
    修改 \codeinline{~/.pip/pip.conf} (没有就创建一个)， 内容如下：
\begin{lstlisting}[style=demostyle,numbers=none]
[global]
index-url = https://pypi.tuna.tsinghua.edu.cn/simple
\end{lstlisting}

\begin{lstlisting}[style=demostyle,numbers=none]
(lora)tony@gusu:~/Lora$ pip config set global.index-url https://pypi.tuna.tsinghua.edu.cn/simple
Writing to /home/bxyang/.config/pip/pip.conf
\end{lstlisting}

    \item Windows
    在用户目录中创建一个 \codeinline{pip} 目录，如 \codeinline{C:\textbackslash{}Users\textbackslash{}tony\textbackslash{}pip}，新建文件 \codeinline{pip.ini}，内容如下：
\begin{lstlisting}[style=demostyle,numbers=none]
[global]
index-url = https://pypi.tuna.tsinghua.edu.cn/simple
\end{lstlisting}

\end{enumerate}


\section{使用}
\subsection{升级}
\begin{lstlisting}[style=demostyle,numbers=none]
# 升级pip
pip install -U pip 
## 如果这个升级命令出现问题 ，可以使用以下命令：
pip install --upgrade pip
\end{lstlisting}

\subsection{搜索包}
\begin{lstlisting}[style=demostyle,numbers=none]
pip search SomePackage
\end{lstlisting}

\subsection{安装包}
\begin{lstlisting}[style=demostyle,numbers=none]
pip install SomePackage              # 最新版本
pip install SomePackage==1.0.4       # 指定版本
pip install 'SomePackage>=1.0.4'     # 最小版本
pip install 'SomePackage>=1.0.4,<2.0' # 指定范围
pip install -i https://pypi.tuna.tsinghua.edu.cn/simple Django # 指定源
\end{lstlisting}

\subsection{查看包}
\begin{lstlisting}[style=demostyle,numbers=none]
pip show SomePackage
# 查看已安装的包
pip list
# 查看指定包的详细信息
pip show --verbose SomePackage
pip show -f SomePackage
\end{lstlisting}

\subsection{升级包}
\begin{lstlisting}[style=demostyle,numbers=none]
#  查看可升级的包
pip list -o
# 升级包
pip install --upgrade SomePackage
\end{lstlisting}

\subsection{卸载包}
\begin{lstlisting}[style=demostyle,numbers=none]
pip uninstall SomePackage
\end{lstlisting}
